\section{The maximum--lexicographic order on ordinal pairs}

\import{order/well-order.tex}
\import{ordinal.tex}

The level of a pair of ordinals is their union, hence their maximum.  Ordering
first by this level and then lexicographically gives a set-bounded version of
the classical well-order of $\ON\times\ON$ used in the proof that an infinite
cardinal is equinumerous with its square.

\begin{definition}\label{ordinal_pair_level}
    $\ordinalPairLevel{p}=\fst{p}\union\snd{p}$.
\end{definition}

\begin{definition}\label{ordinal_pair_precedes}
    $\ordinalPairPrecedes{p}{q}$ iff
        $\ordinalPairLevel{p}\in\ordinalPairLevel{q}$ or
        ($\ordinalPairLevel{p}=\ordinalPairLevel{q}$ and
            ($\fst{p}\in\fst{q}$ or
                ($\fst{p}=\fst{q}$ and $\snd{p}\in\snd{q}$))).
\end{definition}

\begin{definition}\label{ordinal_pair_order}
    $\ordinalPairOrder{\kappa}
        =\{w\in(\kappa\times\kappa)\times(\kappa\times\kappa)\mid
            \ordinalPairPrecedes{\fst{w}}{\snd{w}}\}$.
\end{definition}

\begin{proposition}\label{ordinal_pair_level_is_ordinal}
    Let $\kappa$ be an ordinal.
    Let $p\in\kappa\times\kappa$.
    Then $\ordinalPairLevel{p}$ is an ordinal.
\end{proposition}
\begin{proof}
    $\fst{p},\snd{p}\in\kappa$ by \obv{fst_type,snd_type}.
    Hence $\fst{p}$ and $\snd{p}$ are ordinals
        by \obv{elem_of_ordinal_is_ordinal}.
    Follows by
        \auto{ordinal_pair_level,union_of_two_ordinals_is_ordinal}.
\end{proof}

\begin{proposition}\label{ordinal_pair_coordinates_below_successor_level}
    Let $\kappa$ be an ordinal.
    Let $p\in\kappa\times\kappa$.
    Then $\fst{p},\snd{p}\in\suc{\ordinalPairLevel{p}}$.
\end{proposition}
\begin{proof}
    $\fst{p},\snd{p}\in\kappa$ by \obv{fst_type,snd_type}.
    Hence $\fst{p}$ and $\snd{p}$ are ordinals
        by \obv{elem_of_ordinal_is_ordinal}.
    $\ordinalPairLevel{p}$ is an ordinal
        by \auto{ordinal_pair_level,union_of_two_ordinals_is_ordinal}.
    $\fst{p}\subseteq\fst{p}\union\snd{p}$ and
        $\snd{p}\subseteq\fst{p}\union\snd{p}$
        by \obv{union_upper_left,union_upper_right}.
    Hence $\fst{p}=\ordinalPairLevel{p}$ or
        $\fst{p}\in\ordinalPairLevel{p}$
        by \auto{ordinal_pair_level,ordinal_subseteq_implies_elem_or_eq}.
    $\snd{p}=\ordinalPairLevel{p}$ or
        $\snd{p}\in\ordinalPairLevel{p}$
        by \auto{ordinal_pair_level,ordinal_subseteq_implies_elem_or_eq}.
    Follows by \auto{suc_iff}.
\end{proof}

\begin{proposition}\label{ordinal_pair_order_iff}
    Let $p,q\in\kappa\times\kappa$.
    Then $p\mathrel{\ordinalPairOrder{\kappa}}q$ iff
        $\ordinalPairPrecedes{p}{q}$.
\end{proposition}
\begin{proof}
    Follows by
        \auto{ordinal_pair_order,fst_eq,snd_eq,times_tuple_intro}.
\end{proof}

\begin{proposition}\label{ordinal_pair_order_binary}
    $\ordinalPairOrder{\kappa}$ is a binary relation on
        $\kappa\times\kappa$.
\end{proposition}
\begin{proof}
    Follows by \obv{ordinal_pair_order,subseteq}.
\end{proof}

\begin{proposition}\label{ordinal_pair_precedes_transitive}
    Let $\kappa$ be an ordinal.
    Let $p,q,r\in\kappa\times\kappa$.
    Suppose $\ordinalPairPrecedes{p}{q}$ and
        $\ordinalPairPrecedes{q}{r}$.
    Then $\ordinalPairPrecedes{p}{r}$.
\end{proposition}
\begin{proof}
    $\ordinalPairLevel{r}$ is an ordinal
        by \obv{ordinal_pair_level_is_ordinal}.
    $\fst{p},\snd{p},\fst{q},\snd{q},\fst{r},\snd{r}\in\kappa$
        by \obv{fst_type,snd_type}.
    $\fst{r}$ and $\snd{r}$ are ordinals
        by \obv{elem_of_ordinal_is_ordinal}.
    Follows by \auto{ordinal_pair_precedes,ordinal_transitivity}.
\end{proof}

\begin{proposition}\label{ordinal_pair_precedes_asymmetric}
    Let $\kappa$ be an ordinal.
    Let $p,q\in\kappa\times\kappa$.
    Suppose $\ordinalPairPrecedes{p}{q}$.
    Then not $\ordinalPairPrecedes{q}{p}$.
\end{proposition}
\begin{proof}
    $\fst{p},\snd{p},\fst{q},\snd{q}\in\kappa$
        by \obv{fst_type,snd_type}.
    $\fst{q}$ and $\snd{q}$ are ordinals
        by \obv{elem_of_ordinal_is_ordinal}.
    Follows by \auto{ordinal_pair_precedes,in_asymmetric}.
\end{proof}

\begin{proposition}\label{ordinal_pair_precedes_connex}
    Let $\kappa$ be an ordinal.
    Let $p,q\in\kappa\times\kappa$.
    Suppose $p\neq q$.
    Then $\ordinalPairPrecedes{p}{q}$ or
        $\ordinalPairPrecedes{q}{p}$.
\end{proposition}
\begin{proof}
    $\ordinalPairLevel{p}$ and $\ordinalPairLevel{q}$ are ordinals
        by \obv{ordinal_pair_level_is_ordinal}.
    $\fst{p},\snd{p},\fst{q},\snd{q}\in\kappa$
        by \obv{fst_type,snd_type}.
    Hence $\fst{p}$ and $\snd{p}$ are ordinals
        by \obv{elem_of_ordinal_is_ordinal}.
    $\fst{q}$ and $\snd{q}$ are ordinals
        by \obv{elem_of_ordinal_is_ordinal}.
    $p=(\fst{p},\snd{p})$ and $q=(\fst{q},\snd{q})$
        by \auto{times_elem_is_tuple,fst_eq,snd_eq}.
    Follows by
        \auto{ordinal_pair_precedes,ordinal_elem_connex,pair_eq_iff}.
\end{proof}

\begin{proposition}\label{ordinal_pair_order_transitive}
    Let $\kappa$ be an ordinal.
    Then $\ordinalPairOrder{\kappa}$ is transitive.
\end{proposition}
\begin{proof}
    Show that for all $p,q,r$ such that
        $p\mathrel{\ordinalPairOrder{\kappa}}q$ and
        $q\mathrel{\ordinalPairOrder{\kappa}}r$ we have
        $p\mathrel{\ordinalPairOrder{\kappa}}r$.
    \begin{subproof}
        Fix $p,q,r$.
        Assume $p\mathrel{\ordinalPairOrder{\kappa}}q$ and
            $q\mathrel{\ordinalPairOrder{\kappa}}r$.
        Hence $p,q,r\in\kappa\times\kappa$
            by \auto{ordinal_pair_order,times_tuple_elim}.
        Hence $\ordinalPairPrecedes{p}{q}$ and
            $\ordinalPairPrecedes{q}{r}$
            by \obv{ordinal_pair_order_iff}.
        Hence $\ordinalPairPrecedes{p}{r}$
            by \obv{ordinal_pair_precedes_transitive}.
        Follows by \obv{ordinal_pair_order_iff}.
    \end{subproof}
    Follows by \obv{transitive}.
\end{proof}

\begin{proposition}\label{ordinal_pair_order_asymmetric}
    Let $\kappa$ be an ordinal.
    Then $\ordinalPairOrder{\kappa}$ is asymmetric.
\end{proposition}
\begin{proof}
    Show that for all $p,q$ such that
        $p\mathrel{\ordinalPairOrder{\kappa}}q$ we have
        $q\not\mathrel{\ordinalPairOrder{\kappa}}p$.
    \begin{subproof}
        Fix $p,q$.
        Assume $p\mathrel{\ordinalPairOrder{\kappa}}q$.
        Hence $p,q\in\kappa\times\kappa$
            by \obv{ordinal_pair_order,times_tuple_elim}.
        Hence $\ordinalPairPrecedes{p}{q}$
            by \obv{ordinal_pair_order_iff}.
        Hence not $\ordinalPairPrecedes{q}{p}$
            by \obv{ordinal_pair_precedes_asymmetric}.
        Suppose $q\mathrel{\ordinalPairOrder{\kappa}}p$.
        Hence $\ordinalPairPrecedes{q}{p}$
            by \obv{ordinal_pair_order_iff}.
        Contradiction.
    \end{subproof}
    Follows by \obv{asymmetric}.
\end{proof}

\begin{proposition}\label{ordinal_pair_order_connex}
    Let $\kappa$ be an ordinal.
    Then $\ordinalPairOrder{\kappa}$ is connex on $\kappa\times\kappa$.
\end{proposition}
\begin{proof}
    Follows by
        \auto{connex,ordinal_pair_order_iff,ordinal_pair_precedes_connex}.
\end{proof}

\begin{proposition}\label{ordinal_pair_order_has_least}
    Let $\kappa$ be an ordinal.
    Suppose $X\subseteq\kappa\times\kappa$ and $X\neq\emptyset$.
    Then there exists $p\in X$ such that for every $q\in X$ we have
        $p\mathrel{\ordinalPairOrder{\kappa}}q$ or $p=q$.
\end{proposition}
\begin{proof}
    $X$ is inhabited by \obv{nonempty_inhabited}.
    Let $L=\{\ordinalPairLevel{r}\mid r\in X\}$.
    Take $r_0$ such that $r_0\in X$.
    Hence $\ordinalPairLevel{r_0}\in L$.
    Hence $L$ is inhabited.
    Show that every element of $L$ is an ordinal.
    \begin{subproof}
        Fix $d\in L$.
        Take $r\in X$ such that $d=\ordinalPairLevel{r}$.
        $r\in\kappa\times\kappa$ by \obv{subseteq}.
        Follows by \obv{ordinal_pair_level_is_ordinal}.
    \end{subproof}
    Let $d_0=\inters{L}$.
    $d_0\in L$ by \auto{inters_of_ordinals_elem}.
    Take $r_1\in X$ such that
        $d_0=\ordinalPairLevel{r_1}$.
    Let $X_0=\{r\in X\mid \ordinalPairLevel{r}=d_0\}$.
    $r_1\in X_0$.
    Hence $X_0$ is inhabited.
    Let $F=\{\fst{r}\mid r\in X_0\}$.
    $\fst{r_1}\in F$.
    Hence $F$ is inhabited.
    Show that every element of $F$ is an ordinal.
    \begin{subproof}
        Fix $a\in F$.
        Take $r\in X_0$ such that $a=\fst{r}$.
        Hence $r\in X$.
        Hence $r\in\kappa\times\kappa$ by \obv{subseteq}.
        Hence $\fst{r}\in\kappa$ by \obv{fst_type}.
        Follows by \obv{elem_of_ordinal_is_ordinal}.
    \end{subproof}
    Let $a_0=\inters{F}$.
    $a_0\in F$ by \auto{inters_of_ordinals_elem}.
    Take $r_2\in X_0$ such that $a_0=\fst{r_2}$.
    Let $X_1=\{r\in X_0\mid \fst{r}=a_0\}$.
    $r_2\in X_1$.
    Hence $X_1$ is inhabited.
    Let $G=\{\snd{r}\mid r\in X_1\}$.
    $\snd{r_2}\in G$ by \auto.
    Hence $G$ is inhabited.
    Show that every element of $G$ is an ordinal.
    \begin{subproof}
        Fix $b\in G$.
        Take $r\in X_1$ such that $b=\snd{r}$.
        Hence $r\in X_0$ and $r\in X$ by \auto.
        Hence $r\in\kappa\times\kappa$ by \obv{subseteq}.
        Hence $\snd{r}\in\kappa$ by \obv{snd_type}.
        Follows by \obv{elem_of_ordinal_is_ordinal}.
    \end{subproof}
    Let $b_0=\inters{G}$.
    $b_0\in G$ by \auto{inters_of_ordinals_elem}.
    Take $p\in X_1$ such that $b_0=\snd{p}$.
    Hence $p\in X_0$ and $p\in X$ by \auto.
    Hence $\ordinalPairLevel{p}=d_0$ and $\fst{p}=a_0$ and
        $\snd{p}=b_0$.
    $p\in\kappa\times\kappa$ by \obv{subseteq}.
    Show that for every $q\in X$ we have
        $p\mathrel{\ordinalPairOrder{\kappa}}q$ or $p=q$.
    \begin{subproof}
        Fix $q\in X$.
        Hence $q\in\kappa\times\kappa$ by \obv{subseteq}.
        $\ordinalPairLevel{q}\in L$ by \auto.
        Hence $d_0\subseteq\ordinalPairLevel{q}$
            by \auto{inters_of_ordinals_subseteq}.
        Hence $d_0\in\ordinalPairLevel{q}$ or
            $d_0=\ordinalPairLevel{q}$
            by \obv{ordinal_subseteq_implies_elem_or_eq}.
        \begin{byCase}
            \caseOf{$d_0\in\ordinalPairLevel{q}$.}
                $\ordinalPairPrecedes{p}{q}$
                    by \auto{ordinal_pair_precedes}.
                Follows by \obv{ordinal_pair_order_iff}.
            \caseOf{$d_0=\ordinalPairLevel{q}$.}
                Hence $q\in X_0$.
                Hence $\fst{q}\in F$ by \auto.
                Hence $a_0\subseteq\fst{q}$
                    by \auto{inters_of_ordinals_subseteq}.
                Hence $a_0\in\fst{q}$ or $a_0=\fst{q}$
                    by \obv{ordinal_subseteq_implies_elem_or_eq}.
                \begin{byCase}
                    \caseOf{$a_0\in\fst{q}$.}
                        $\ordinalPairPrecedes{p}{q}$
                            by \auto{ordinal_pair_precedes}.
                        Follows by \obv{ordinal_pair_order_iff}.
                    \caseOf{$a_0=\fst{q}$.}
                        Hence $q\in X_1$.
                        Hence $\snd{q}\in G$ by \auto.
                        Hence $b_0\subseteq\snd{q}$
                            by \auto{inters_of_ordinals_subseteq}.
                        Hence $b_0\in\snd{q}$ or $b_0=\snd{q}$
                            by \obv{ordinal_subseteq_implies_elem_or_eq}.
                        \begin{byCase}
                            \caseOf{$b_0\in\snd{q}$.}
                                $\ordinalPairPrecedes{p}{q}$
                                    by \auto{ordinal_pair_precedes}.
                                Follows by \obv{ordinal_pair_order_iff}.
                            \caseOf{$b_0=\snd{q}$.}
                                $p=(\fst{p},\snd{p})$ and
                                    $q=(\fst{q},\snd{q})$ by
                                    \auto{times_elem_is_tuple,fst_eq,snd_eq}.
                                Follows by \obv{pair_eq_iff}.
                        \end{byCase}
                \end{byCase}
        \end{byCase}
    \end{subproof}
    Follows by \auto.
\end{proof}

\begin{theorem}\label{ordinal_pair_order_is_strict_well_order}
    Let $\kappa$ be an ordinal.
    Then $\ordinalPairOrder{\kappa}$ is a strict well order on
        $\kappa\times\kappa$.
\end{theorem}
\begin{proof}
    Follows by \auto{strict_well_order_iff,ordinal_pair_order_binary,ordinal_pair_order_transitive,ordinal_pair_order_asymmetric,ordinal_pair_order_connex,ordinal_pair_order_has_least}.
\end{proof}

\begin{proposition}\label{ordinal_pair_level_belongs}
    Let $\kappa$ be an ordinal.
    Let $p\in\kappa\times\kappa$.
    Then $\ordinalPairLevel{p}\in\kappa$.
\end{proposition}
\begin{proof}
    $\fst{p},\snd{p}\in\kappa$ by \obv{fst_type,snd_type}.
    Hence $\fst{p}$ and $\snd{p}$ are ordinals
        by \obv{elem_of_ordinal_is_ordinal}.
    $\fst{p}\subseteq\snd{p}$ or $\snd{p}\subseteq\fst{p}$
        by \obv{ordinal_subseteq_or_subseteq}.
    Hence $\fst{p}\union\snd{p}=\snd{p}$ or
        $\fst{p}\union\snd{p}=\fst{p}$
        by \obv{union_absorb_subseteq_left,union_absorb_subseteq_right}.
    Follows by \auto{ordinal_pair_level}.
\end{proof}

\begin{proposition}\label{ordinal_pair_predecessor_bound}
    Let $\kappa$ be an ordinal.
    Let $p,q\in\kappa\times\kappa$.
    Suppose $p\mathrel{\ordinalPairOrder{\kappa}}q$.
    Then $p\in
        \suc{\ordinalPairLevel{q}}\times\suc{\ordinalPairLevel{q}}$.
\end{proposition}
\begin{proof}
    $\ordinalPairLevel{p}$ and $\ordinalPairLevel{q}$ are ordinals
        by \obv{ordinal_pair_level_is_ordinal}.
    $\ordinalPairPrecedes{p}{q}$ by \obv{ordinal_pair_order_iff}.
    Hence $\ordinalPairLevel{p}\in\ordinalPairLevel{q}$ or
        $\ordinalPairLevel{p}=\ordinalPairLevel{q}$
        by \auto{ordinal_pair_precedes}.
    Hence $\ordinalPairLevel{p}\subseteq\ordinalPairLevel{q}$
        by \obv{ordinal_elem_implies_subseteq,subseteq_refl}.
    Hence $\suc{\ordinalPairLevel{p}}
        \subseteq\suc{\ordinalPairLevel{q}}$
        by \auto{subseteq,suc_iff,ordinal_subseteq_implies_elem_or_eq}.
    $\fst{p},\snd{p}\in\suc{\ordinalPairLevel{p}}$
        by \obv{ordinal_pair_coordinates_below_successor_level}.
    Hence $\fst{p},\snd{p}\in\suc{\ordinalPairLevel{q}}$
        by \obv{subseteq}.
    $p=(\fst{p},\snd{p})$
        by \auto{times_elem_is_tuple,fst_eq,snd_eq}.
    Follows by \auto{times_tuple_intro}.
\end{proof}

\begin{corollary}\label{ordinal_pair_predecessor_bound_in_limit}
    Let $\kappa$ be a limit ordinal.
    Let $p,q\in\kappa\times\kappa$.
    Suppose $p\mathrel{\ordinalPairOrder{\kappa}}q$.
    Then there exists $\sigma\in\kappa$ such that
        $p\in\sigma\times\sigma$.
\end{corollary}
\begin{proof}
    $\ordinalPairLevel{q}\in\kappa$
        by \obv{ordinal_pair_level_belongs}.
    Hence $\suc{\ordinalPairLevel{q}}\in\kappa$
        by \obv{suc_elem_limitordinal}.
    $p\in\suc{\ordinalPairLevel{q}}
        \times\suc{\ordinalPairLevel{q}}$
        by \obv{ordinal_pair_predecessor_bound}.
\end{proof}
